% =====================================================================
%  DASC7606 Assignment 1 -- Final Report (PDF, up to 2 pages)
%  Compile:  pdflatex report.tex  (or use overleaf.com)
% =====================================================================
\documentclass[10pt,conference]{article}
\usepackage[utf8]{inputenc}
\usepackage{graphicx}
\usepackage{booktabs}
\usepackage{amsmath}
\usepackage{hyperref}
\usepackage[margin=0.85in]{geometry}
\usepackage{caption}
\captionsetup{font=small,labelfont=bf}

\title{DASC7606 Assignment 1: Image Classification on CIFAR-10}
\author{Your Name \\ Your UID: \texttt{3036xxxxxx}}
\date{\today}

\begin{document}
\maketitle

% -------------------------------------------------------------------
\section{Introduction}
% -------------------------------------------------------------------
Image classification is a foundational task in computer vision.
In this assignment we train a convolutional neural network (CNN)
to recognise the ten object categories of the
\textit{CIFAR-10} dataset~\cite{krizhevsky2009cifar}.
The dataset contains 60{,}000 colour images of size $32 \times 32$
(50{,}000 for training, 10{,}000 for testing) and serves as a long-standing
benchmark for new architectures and training techniques.

The first deep CNN to dominate CIFAR-10 was VGG~\cite{simonyan2015vgg},
which stacks small $3\times3$ convolutions and reaches $\sim$92\% test
accuracy with standard augmentation. Subsequent work showed that
performance can be pushed above 95\% by combining strong augmentation
(\textit{Mixup}~\cite{zhang2018mixup}, \textit{CutMix}~\cite{yun2019cutmix},
\textit{AutoAugment}~\cite{cubuk2019autoaugment}, \textit{RandomErasing})
with regularisation techniques such as \textit{Batch Normalisation},
\textit{Dropout} and \textit{label smoothing}, plus a
\textit{cosine learning-rate schedule}~\cite{loshchilov2017sgdr}.

\textbf{Contributions.}  We (i) reproduce a 3-layer baseline;
(ii) deepen it to a 5-layer network with BatchNorm and Dropout;
(iii) design an open-ended VGG-16 style model and train it with the
strong regularisation recipe above, reaching our best test accuracy.

% -------------------------------------------------------------------
\section{Methods}
% -------------------------------------------------------------------
\paragraph{Exercise 1 -- Baseline.}  Three conv layers
($3 \to 32 \to 64 \to 128$) with ReLU and $2\times2$ max-pool,
followed by two fully-connected layers
($2048 \to 256 \to 10$).  No regularisation; trained with plain SGD
for 20 epochs.

\paragraph{Exercise 2 -- Improved network.}  Five conv layers arranged
in three conv-blocks ($3 \to 64 \to 128 \to 256$), each block
followed by BatchNorm, ReLU and max-pool.  Classifier is
$4096 \to 512 \to 256 \to 10$ with Dropout($0.3$).

\paragraph{Exercise 3 -- Open-ended design (\emph{submitted model}).}
We adopt a VGG-16 style backbone
(\texttt{cfg~=~[64,64,M, 128,128,M, 256,256,256,M, 512,512,512,M, 512,512,512,M]})
with BatchNorm after every convolution, AdaptiveAvgPool before the
classifier and Dropout($0.3$) in the head ($\sim$14.7M parameters).
Training regularisation:
\begin{itemize}\setlength\itemsep{0pt}
  \item Augmentation: \texttt{RandomCrop(32, pad=4)} +
        \texttt{RandomHorizontalFlip} + \texttt{AutoAugment(CIFAR10)} +
        \texttt{RandomErasing($p{=}0.25$)}.
  \item Mixup or CutMix ($\alpha{=}1.0$, each batch chosen with
        probability 0.5).
  \item Label smoothing $\varepsilon{=}0.1$.
  \item Optimiser: SGD with Nesterov momentum $0.9$, weight decay
        $5\!\times\!10^{-4}$, initial LR $0.1$.
  \item LR schedule: cosine annealing to 0 over 100 epochs,
        batch size 128.
\end{itemize}

\begin{table}[h]
\centering
\small
\begin{tabular}{l r r}
\toprule
\textbf{Backbone} & \textbf{Params (M)} & \textbf{FLOPs (G)} \\
\midrule
BaselineCNN (Exp1) &  0.62 &  XX.X \\
ImprovedCNN (Exp2) &  2.79 &  XX.X \\
\textbf{VGG-16 (Exp3)}  & \textbf{14.73} & \textbf{XX.X} \\
\bottomrule
\end{tabular}
\caption{Parameter and FLOPs budget of the three models (fill in FLOPs from \texttt{model\_stats.py}).}
\label{tab:params}
\end{table}

\paragraph{Hyper-parameter selection.}  We performed a short
learning-rate sensitivity sweep (initial LR $\in \{0.01, 0.05, 0.1,
0.5\}$, see \texttt{lr\_sweep.py}).  LR$=$0.1 with a cosine annealing
schedule gave the best validation accuracy and was used for the final
model.

% -------------------------------------------------------------------
\section{Experiments \& Analysis}
% -------------------------------------------------------------------
\textbf{Setup.}  We split the official 50k training images into a
45k/5k train/validation split and report the test-set accuracy of
the checkpoint that achieved the highest validation accuracy.
All experiments are run on a single GPU of the HKU GPU Farm.

\subsection*{Results}
\begin{table}[h]
\centering
\small
\begin{tabular}{l l c c c}
\toprule
\textbf{Experiment} & \textbf{Backbone} &
\textbf{Test Acc (\%)} &
\textbf{Epochs} & \textbf{Augmentation} \\
\midrule
Exp1 Baseline    & 3 conv + 2 FC          & XX.X &  20 & None \\
Exp2 Improved    & 5 conv + BN/Dropout    & XX.X &  50 & Basic flip/crop \\
Exp3 \textbf{Ours} & VGG-16 + BN/Dropout   & \textbf{XX.X} & 100 & Full recipe \\
\bottomrule
\end{tabular}
\caption{Comparison of the three experiments.}
\label{tab:results}
\end{table}

\subsection*{Training curves}
\begin{figure}[h]
\centering
\begin{minipage}{0.49\linewidth}
  \centering
  \includegraphics[width=\linewidth]{training_curves.png}
  \caption{Training/validation loss and accuracy of the Exp3 model.}
  \label{fig:curve}
\end{minipage}\hfill
\begin{minipage}{0.49\linewidth}
  \centering
  \includegraphics[width=\linewidth]{augmentation_grid.png}
  \caption{Raw (top) and strongly-augmented (bottom) CIFAR-10 samples.}
  \label{fig:aug}
\end{minipage}
\end{figure}

\begin{figure}[h]
\centering
\begin{minipage}{0.49\linewidth}
  \centering
  \includegraphics[width=\linewidth]{confusion_matrix.png}
  \caption{Normalised 10$\times$10 confusion matrix on the test set.}
  \label{fig:cm}
\end{minipage}\hfill
\begin{minipage}{0.49\linewidth}
  \centering
  \includegraphics[width=\linewidth]{per_class_accuracy.png}
  \caption{Per-class test accuracy.}
  \label{fig:pca}
\end{minipage}
\end{figure}

\begin{figure}[h]
\centering
\begin{minipage}{0.49\linewidth}
  \centering
  \includegraphics[width=\linewidth]{misclassified.png}
  \caption{Most confident misclassifications.}
  \label{fig:miss}
\end{minipage}\hfill
\begin{minipage}{0.49\linewidth}
  \centering
  \includegraphics[width=\linewidth]{tsne.png}
  \caption{t-SNE of the penultimate features.}
  \label{fig:tsne}
\end{minipage}
\end{figure}

\subsection*{Analysis}
\begin{itemize}\setlength\itemsep{2pt}
  \item \textbf{Depth vs.\ width.}  Going from Exp1 to Exp2
        (3$\to$5 conv blocks + BN) gives a $+$\textit{XX.X}\% absolute
        test-accuracy gain, mostly because BatchNorm stabilises
        optimisation and allows deeper networks.
  \item \textbf{Regularisation is decisive.}  Switching the Exp2 model
        to the Exp3 training recipe (AutoAugment + Mixup/CutMix +
        label smoothing + cosine schedule) further improves
        accuracy by another $+$\textit{XX.X}\%, with the largest
        contribution coming from the augmentation pipeline.
  \item \textbf{Per-class error.}  The confusion matrix (see
        supplementary notebook) shows that the model most often
        confuses \emph{cat} with \emph{dog} and \emph{deer} with
        \emph{horse} -- two semantically close class pairs.
  \item \textbf{Ablation.}  We disable each component of the
        training recipe in turn (see Table~\ref{tab:ablation});
        removing AutoAugment and Mixup/CutMix each cause a drop
        of $1$--$2\%$, confirming that both are needed.
\end{itemize}

\begin{table}[h]
\centering
\small
\begin{tabular}{l c c}
\toprule
\textbf{Configuration} & \textbf{Val Acc (\%)} & \textbf{Test Acc (\%)} \\
\midrule
Full recipe (Exp3)               & XX.X & \textbf{XX.X} \\
\quad $-$ AutoAugment            & XX.X & XX.X \\
\quad $-$ Mixup / CutMix         & XX.X & XX.X \\
\quad $-$ RandomErasing          & XX.X & XX.X \\
\quad $-$ Label smoothing        & XX.X & XX.X \\
\quad $-$ BatchNorm              & XX.X & XX.X \\
\quad $-$ Dropout                & XX.X & XX.X \\
\quad Cosine $\to$ Step LR       & XX.X & XX.X \\
\bottomrule
\end{tabular}
\caption{Ablation study on the Exp3 model (values filled in from \texttt{ablation\_results.csv}).}
\label{tab:ablation}
\end{table}

\subsection*{Robustness}
Figure~\ref{fig:robust} reports the test accuracy as a function of
Gaussian noise $\sigma$ and cutout patch size.  Even at $\sigma=0.3$
the model retains $>$XX\% accuracy, and it stays $>$XX\% with a
$16\!\times\!16$ occluding patch.  This is a direct consequence of
the strong augmentation pipeline used at training time.

\begin{figure}[h]
\centering
\includegraphics[width=0.7\linewidth]{robustness.png}
\caption{Robustness to additive Gaussian noise (left) and cutout
  occlusion (right).}
\label{fig:robust}
\end{figure}

\subsection*{Conclusion}
We presented a three-stage study on CIFAR-10, showing that a VGG-16
style backbone, when combined with strong augmentation and modern
regularisation, reaches a high test accuracy without any pre-training.
Future work could explore attention mechanisms (SE-Net) or lighter
architectures (ResNet-18) to improve the accuracy / parameter trade-off.

% -------------------------------------------------------------------
\begin{thebibliography}{9}
\bibitem{krizhevsky2009cifar} A. Krizhevsky, \emph{Learning Multiple
  Layers of Features from Tiny Images}, 2009.
\bibitem{simonyan2015vgg}    K. Simonyan \& A. Zisserman,
  \emph{Very Deep Convolutional Networks for Large-Scale Image
  Recognition}, ICLR 2015.
\bibitem{zhang2018mixup}     H. Zhang et al., \emph{mixup: Beyond
  Empirical Risk Minimization}, ICLR 2018.
\bibitem{yun2019cutmix}      S. Yun et al., \emph{CutMix: Regularization
  Strategy to Train Strong Classifiers with Localizable Features},
  ICCV 2019.
\bibitem{cubuk2019autoaugment} E. D. Cubuk et al.,
  \emph{AutoAugment: Learning Augmentation Strategies From Data},
  CVPR 2019.
\bibitem{loshchilov2017sgdr} I. Loshchilov \& F. Hutter,
  \emph{SGDR: Stochastic Gradient Descent with Warm Restarts},
  ICLR 2017.
\end{thebibliography}

\end{document}
